\section{Experimental Protocol}

\implemented{} The experiment infrastructure identifies formal comparisons by
immutable IDs and links them to a
versioned configuration, Git state, dependency versions, vehicle and dock
fingerprints, seeds, output paths, validity status, and interpretation notes.
MuJoCo runs will additionally store the complete resolved tuning values because
those values are not part of the version-1 vehicle fingerprint.
Large checkpoints and trajectories will remain outside the tracked manuscript,
while compact metrics, summaries, and reproducible figures will be retained.
The version-1 suite assigns training seeds 0--2, validation episodes beginning
at seed 100000, and a disjoint 100-episode final test range beginning at seed
200000. These are split definitions, not completed experiments. Validation may
guide tuning; the held-out range must remain untouched until the condition is
frozen.

\implemented{} Reward tuning can be executed as a controlled study rather than
as unrelated training commands. Conditions are constructed from one base
mixed-curve scenario and may change only declared reward coefficients. Every
condition is trained on the same selected seed set and validated on the same
physical starts, geometries, and sampled curve templates. Because a reward
change alters the scenario fingerprint, both full episode fingerprints and
reward-independent physical-input fingerprints are retained. Validation
results may select a condition but are not held-out test evidence.

The primary comparison will evaluate PPO and NMPC on identical held-out starts,
target docks, geometry, plant limits, state availability, and safety definitions.
The implemented episode resolver freezes each sampled geometry, articulated
start state, and reference path under a content fingerprint, and the comparison
runner replays that episode set for every selected method and plant. Thus the
pairing mechanism exists even though the held-out study has not yet been run.
Methodological asymmetries will be stated explicitly: NMPC receives an explicit
prediction model, while PPO receives training experience and a selected
observation representation. Offline training/tuning cost will be separated
from online control computation.

Planned metrics are:
\begin{itemize}
  \item docking success, collision, and jackknife rates;
  \item final longitudinal, lateral, and trailer-heading error;
  \item maximum articulation angle;
  \item maneuver time and path length;
  \item steering-rate magnitude and smoothness;
  \item controller computation time and numerical failures; and
  \item safety-supervisor intervention count.
\end{itemize}

After nominal performance is established, candidate robustness studies include
held-out start poses, hitch-to-axle distances, bay dimensions, clearance, state
noise, and model mismatch. Geometry variation alone will be described as
sensitivity, robustness, or generalization; the term \emph{ablation} will be
reserved for removing or modifying a defined component.

Transfer evaluation will proceed only after 2D/MuJoCo parity checks establish
matching frames, signs, dimensions, command timing, semantic sites, and event
definitions. Zero-shot transfer, domain-randomized transfer, and MuJoCo
fine-tuning will be reported as separate conditions if they are implemented.
